\documentclass[11pt]{article}
\usepackage[margin=1in]{geometry}
\usepackage{fontspec}
\usepackage{setspace}
\usepackage{hyperref}
\hypersetup{colorlinks=true, allcolors=black}
\setlength{\parindent}{0pt}
\setlength{\parskip}{0.55em}
\setstretch{1.12}
\pagestyle{empty}
\setmainfont{Times New Roman}

\begin{document}

\begin{center}
{\Large Teaching Philosophy}\\[0.35em]
Ke Cheng HSU\\
International Business Studies, Temple University, Japan Campus
\end{center}

\noindent\textbf{Clarity first.} I begin with explanation before practice. For each concept I break the reasoning into small, ordered steps and work through a live example before asking students to try a short in-class exercise on their own. An explanation students can follow, immediately followed by a demonstration, lowers the barrier they feel when they meet new quantitative or technical material. At HKUST I used the same approach as a teaching assistant from 2020 to 2025 for undergraduate macroeconomics, the economics of uncertainty and information, and Ph.D.\ microeconomic theory, with an average evaluation of 4.72 out of 5. In 2026 I designed and lectured the mathematics camp for the MSc in Economics for students from a non-math background, building each tool step by step and receiving 4.53 out of 5.

\noindent\textbf{Business problems first.} Students engage more when a technique answers a question that matters. I open each topic with a real organizational or market problem. How should an online retailer set the terms of its return policy so that it encourages efficient purchases without inviting excessive returns? How can transaction data reveal distinct customer segments? How does a platform's design shape what users buy and how often they return? Anchoring the method in questions like these shows students that analytics and programming are tools for understanding and improving real business decisions, not ends in themselves.

\noindent\textbf{Interactive technology.} I use technology to make the class more participatory, not only more efficient. I run short live polls, word clouds, and real-time quizzes to turn listening into doing. When the class's answers appear together on screen, I can see where understanding is thin and adjust the pace immediately, and students see that others share the same questions. As Teaching Assistant Coordinator for the HKUST Business School from 2025 to 2026, I helped postgraduate instructors build these interactive checkpoints into their teaching materials and mock lectures.

\noindent\textbf{Inclusive, active classroom.} I want every question to be welcome, especially in a class where students come from different countries and majors. After the intuition is clear, I ask students to work through the reasoning and tell me where it stops being clear; one student's question often helps several others and shows me where to explain more simply. Teaching is most rewarding when a student who thought the material was out of reach sees why an outcome happens, and finding that clearer explanation is the part of teaching I enjoy most.

\end{document}
